\sectioncentered*{Заключение}
\addcontentsline{toc}{section}{Заключение}

В ходе выполнения дипломного проекта была успешно спроектирована и реализована интеллектуальная система мониторинга рисков и обеспечения безопасности путешественников, объединяющая в себе передовые методы интеллектуального анализа текстов, геоинформационные технологии и принципы построения высоконагруженных распределенных систем. Работа над проектом охватила полный жизненный цикл создания программного продукта — от глубокого исследования предметной области и анализа существующих рыночных решений до физической реализации микросервисной архитектуры и проведения комплексного тестирования на реальных данных о глобальных инцидентах.

В рамках выполнения поставленных задач были достигнуты следующие результаты:

\begin{itemize}
    \item изучены современные стандарты управления рисками в путешествиях, включая международный стандарт ISO 31030, что позволило сформулировать требования к системе на уровне мировых практик обеспечения безопасности;
    \item исследованы существующие на мировом рынке коммерческие аналоги, такие как International SOS, Safeture и Crisis24, в результате чего была выявлена и обоснована ниша для создания доступного и технологически независимого решения для частных лиц;
    \item сформулированы функциональные и нефункциональные требования к платформе, обеспечивающие высокую скорость реакции системы на возникающие угрозы и возможность глобального масштабирования;
    \item показано преимущество использования событийно-ориентированной микросервисной архитектуры перед традиционными монолитными подходами в контексте обработки непредсказуемых потоков новостных данных;
    \item разработана и внедрена многослойная таксономия рисков, охватывающая геополитические, природные, медицинские и техногенные угрозы, что стало основой для классификации инцидентов нейронными сетями;
    \item предложена и реализована концепция «динамических зон влияния» инцидентов, заменяющая точечное отображение событий на карте построением полигонов риска с учетом их интенсивности;
    \item подготовлены и развернуты микросервисы на языках Python и Java, обеспечивающие полную автоматизацию цикла «сбор — анализ — уведомление» без участия человеческого фактора;
    \item изготовлена и испытана система интеллектуального анализа текстов на базе модели DeBERTa-v3, реализующая метод Zero-shot классификации для распознавания новых видов угроз в режиме реального времени;
    \item испытана и подтверждена эффективность использования локальных больших языковых моделей (LLM) семейства Qwen для автоматической генерации кратких рефератов новостей, адаптированных под формат мобильных уведомлений;
    \item разработана структура базы данных на основе PostgreSQL и PostGIS, оптимизированная для выполнения высокоскоростных пространственных запросов на пересечение маршрутов пользователей с зонами риска;
    \item подготовлены инструменты визуализации на базе React и MapLibre GL JS, обеспечивающие плавное отображение глобальной карты инцидентов с использованием возможностей аппаратного ускорения графики;
    \item подтверждено достижение целевых показателей быстродействия системы, при которых время от момента фиксации события в первичном источнике до доставки уведомления пользователю не превышает заданных лимитов.
\end{itemize}

В процессе аналитического исследования была детально изучена база данных GDELT как основной источник информации открытой разведки (OSINT), что позволило выявить закономерности в структуре новостных потоков и настроить алгоритмы фильтрации шума. Исследованы возможности интеграции различных аналитических моделей, в результате чего был сделан обоснованный выбор в пользу гибридного подхода, сочетающего точность классических трансформеров и гибкость генеративных моделей.

На этапе проектирования была разработана детальная архитектурная схема, включающая модули Fetcher, Analyzer и Core Service, взаимодействие между которыми организовано через промышленный брокер сообщений RabbitMQ. Данное решение обеспечило необходимую устойчивость системы к пиковым нагрузкам, возникающим в моменты крупных мировых катастроф, когда объем входящих сообщений возрастает в десятки раз. Сформулированы принципы асинхронного геофенсинга, которые позволяют системе одновременно отслеживать тысячи перемещающихся пользователей без деградации производительности серверной части.

Программная реализация проекта подтвердила правильность выбранного технологического стека. Разработанная на Java серверная часть продемонстрировала высокую надежность в управлении пространственными данными через Hibernate Spatial, а Python-сервисы машинного обучения показали отличные результаты в задачах семантического разбора зашумленных текстов из социальных сетей и новостных лент. Показано, что использование Docker-контейнеризации существенно упрощает процесс развертывания системы в облачных инфраструктурах и обеспечивает изоляцию библиотек машинного обучения от бизнес-логики приложения.

Особое внимание в заключительной фазе проекта было уделено вопросам практической применимости и экономической эффективности. Разработанная система обладает потенциалом для значительного снижения затрат на мониторинг безопасности в туристическом секторе за счет полной автоматизации аналитической работы. Предложенная модель открытого API позволяет интегрировать сервис в существующие корпоративные системы управления командировками, повышая уровень защищенности сотрудников.

Практическая значимость работы заключается в создании работающего прототипа интеллектуальной системы, способной в автономном режиме обеспечивать принцип «должной осмотрительности» (Duty of Care). Реализованные механизмы анализа маршрутов и превентивного оповещения позволяют минимизировать риски попадания путешественников в зоны конфликтов, стихийных бедствий или эпидемиологических вспышек.

В качестве направлений для дальнейшего развития проекта могут быть рассмотрены следующие аспекты:
\begin{itemize}
    \item расширение перечня источников данных за счет интеграции специализированных API мониторинга авиационного и железнодорожного трафика для более точного прогнозирования задержек;
    \item внедрение прогностических моделей на базе графовых нейронных сетей для предсказания динамики распространения протестных настроений или эпидемий в соседние регионы;
    \item реализация механизмов двухсторонней экстренной связи (SOS) с интеграцией локальных служб спасения через протоколы спутниковой передачи данных;
    \item адаптация интерфейса для работы в условиях крайне низкой пропускной способности сети с использованием прогрессивных веб-приложений (PWA);
    \item разработка алгоритмов персонализированной оценки рисков на основе медицинских показателей пользователя и его предыдущего опыта поездок;
    \item интеграция с системами блокчейн для обеспечения неизменяемости истории уведомлений в целях последующего юридического аудита действий служб безопасности.
\end{itemize}

Подводя итого, следует отметить, что все цели и задачи, поставленные в начале дипломного проектирования, выполнены в полном объеме. Созданная интеллектуальная система мониторинга представляет собой законченное техническое решение, объединяющее наукоемкие алгоритмы искусственного интеллекта с практической потребностью современного общества в обеспечении безопасности и информированности в условиях глобальной нестабильности. Полученные результаты могут быть использованы как база для разработки коммерческого продукта или в качестве фундамента для дальнейших научно-исследовательских изысканий в области прикладного анализа больших данных и ГИС-технологий.
